% !TeX root = ../main.tex

\chapter{相关工作}

本章综述与大模型资源效率优化相关的研究工作，涵盖大模型推理系统、分布式训练并行技术以及自动并行化方法三个主要方向，分析现有工作的局限性与本文工作的创新点。

\section{大模型推理系统研究}

大模型推理系统的核心目标是在满足服务质量约束的前提下，最大化系统吞吐量或最小化资源消耗。本节从推理系统架构、能效优化方法和Prefill-Decode分离架构三个方面综述相关研究。

\subsection{推理系统架构与优化技术}

随着大模型在对话系统、代码生成和智能体应用中的广泛部署，推理系统的效率优化成为研究热点。现有工作主要从以下几个方向展开：

\textbf{批处理策略优化。}传统的静态批处理难以适应生成式模型输出长度不确定的特点。Orca\cite{yu2022orca}提出了迭代级调度（iteration-level scheduling）和选择性批处理（selective batching），允许请求在不同迭代步骤动态加入或退出批处理，显著提升了GPU利用率。SGLang\cite{zheng2024sglang}进一步提出了radixAttention技术，通过基数树结构高效管理KV缓存，支持前缀缓存和共享。vLLM\cite{kwon2023efficient}引入了PagedAttention机制，借鉴操作系统的虚拟内存管理思想，将KV缓存划分为固定大小的页面，实现了高效的内存分配和共享。

\textbf{内存管理优化。}大模型推理的显存消耗主要来自模型参数和KV缓存。FlexGen\cite{sheng2023flexgen}提出了CPU卸载策略，将无法装入GPU显存的数据流式传输到CPU内存甚至SSD，通过细粒度的调度隐藏I/O延迟。vAttention\cite{prabhu2025vattention}针对连续批处理场景下的注意力计算进行优化，通过虚拟化技术管理稀疏的注意力矩阵。LightSeq\cite{wang2021lightseq}从GPU内核层面优化Transformer推理，包括自定义CUDA内核和量化技术。

\textbf{投机解码。}投机解码（Speculative Decoding）\cite{miao2024specinfer,stern2018blockwise}通过使用小型草稿模型预测多个token，然后由目标模型并行验证，在不改变输出分布的前提下加速生成过程。SpecInfer\cite{miao2024specinfer}提出了树形投机解码结构，支持多个草稿模型的组合。这一技术对降低端到端延迟有显著效果，但对能效优化的影响尚未被深入探索。

\textbf{调度策略优化。}FastServe\cite{wu2023fast}针对多租户推理场景，提出了抢占式调度策略以降低作业完成时间。AlpaServe\cite{li2023alpaserve}研究了多模型服务场景下的资源分配和请求路由问题。这些工作主要关注延迟和吞吐量优化，较少涉及能效考量。

\subsection{推理能效优化方法}

随着大模型推理能耗问题日益突出，能效优化开始受到关注。现有工作主要从以下几个角度展开：

\textbf{频率调节与功率管理。}DynamoLLM\cite{stojkovic2025dynamollm}首次系统研究了GPU频率控制在大模型推理中的应用，提出基于请求特征（预测的输入/输出token数）的分层频率调节策略。然而，该方法采用粗粒度的频率调整（以分钟为时间尺度），难以适应推理负载的快速波动。$\mu$-Serve\cite{qiu2024power}研究了多模型共服务场景下的功率优化，通过联合调度多个ML模型实现更好的功率效率。ThrottLL'eM\cite{kakolyris2025throttll}通过预测KV缓存使用量和batch size来动态调整频率和实例规模。

\textbf{碳排放与生命周期优化。}EcoServe\cite{li2025ecoserve}从生命周期角度研究推理系统的碳排放问题，提出基于离线/在线推断、负载需求和模型规模的GPU/CPU分配策略。TAPAS\cite{stojkovic2025tapas}关注数据中心的热管理，通过将请求路由到特定机架行/列来利用热和功率的松弛空间，避免高功率热点。Heron\cite{reddy2025ai}提出将GPU部署在靠近可再生能源的位置，并通过负载调整提高能源效率。

\textbf{训练阶段能效优化。}Perseus\cite{zhou2023perseus}研究了分布式训练中的能效问题，发现梯度同步造成的"能效散漫者"（energy straggler）现象：等待梯度聚合的空闲时段使GPU在满功耗状态下空转，浪费大量能耗。Perseus通过识别关键路径外的冗余计算并降低相应算子的功率来节能，在不牺牲训练效果的前提下实现了显著的能效提升。然而，Perseus主要针对数据并行的全局同步点，未深入探索流水线并行中各阶段计算强度差异带来的频率优化机会。

上述工作存在以下局限性：首先，它们大多采用粗粒度的频率控制，无法适应推理负载在毫秒级时间尺度上的变化；其次，它们未区分prefill和decode两个阶段的异质性，采用统一的控制策略；第三，它们未深入分析大模型推理能耗-延迟关系的非线性特征和硬件架构粒度效应。本文的GreenServe系统正是针对这些局限而设计，首次实现了迭代级响应式频率控制和阶段感知的能效优化。

\subsection{Prefill-Decode分离架构}

Prefill阶段和Decode阶段具有截然不同的计算特征：Prefill阶段处理完整输入序列，计算密集且延迟可预测；Decode阶段逐token生成，内存带宽受限且延迟与历史token数相关。这种差异催生了Prefill-Decode分离架构的研究。

SplitWise\cite{patel2024splitwise}首次系统分析了分离架构的优势，指出分离可以避免阶段间干扰、实现异构硬件分配和差异化能耗管理。DistServe\cite{zhong2024distserve}进一步优化了分离架构下的资源分配和请求调度，显著提升了TTFT和ITL的SLO达成率。TetriInfer\cite{hu2024inference}提出基于预测的动态分离策略，根据实时负载动态调整prefill和decode实例的比例。Llumnix\cite{sun2024llumnix}研究了分离架构下的请求迁移和负载均衡问题。

主流推理框架如vLLM\cite{kwon2023efficient}和SGLang\cite{zheng2024sglang}也已原生支持分离架构。然而，这些工作主要关注延迟和吞吐量优化，鲜有将分离架构作为能效优化平台的研究。本文的GreenServe系统首次将分离架构与能效优化深度结合，探索了分离架构下阶段感知频率控制和请求路由的新机会。

\section{分布式训练并行技术研究}

随着大模型参数规模和训练数据量的爆炸式增长，分布式训练已成为必要手段。本节综述主流的并行技术及其演进，重点分析它们在长序列场景下的局限性。

\subsection{数据并行及其变体}

数据并行（Data Parallelism, DP）是最直观的并行策略，将数据集划分到多个设备，每个设备持有完整的模型副本，通过周期性的梯度同步（all-reduce）保持一致性。DP的优点是概念简单、易于实现，缺点是显存占用与设备数无关，难以训练超大模型。

为了解决显存瓶颈，ZeRO（Zero Redundancy Optimizer）系列\cite{rajbhandari2020zero,zhao2023pytorch}引入了状态分片思想。ZeRO-1将优化器状态分片，ZeRO-2进一步分片梯度，ZeRO-3（即FSDP）将模型参数也进行分片。通过all-gather在计算时临时收集所需参数，计算完成后立即释放，ZeRO-3将显存占用从与模型规模线性相关降低到与设备数反比相关。

然而，ZeRO系列的代价是显著增加的通信量。每次前向和反向传播都需要多次all-gather和reduce-scatter操作，且通信量与模型规模成正比。在互联带宽受限的环境中，通信开销可能成为瓶颈。此外，ZeRO的通信模式是集体通信（collective communication），对网络拓扑和带宽有较高要求。

\subsection{模型并行技术}

模型并行（Model Parallelism）将模型本身切分到多个设备，主要包括张量并行和流水线并行两种形式。

\textbf{张量并行。}张量并行（Tensor Parallelism, TP）\cite{shoeybi2019megatron}将单个算子（如矩阵乘法）切分到多个设备并行执行。Megatron-LM\cite{shoeybi2019megatron}系统性地研究了Transformer模型中各层的TP切分策略，包括注意力层和前馈网络层的行/列切分组合。TP的优点是可以大幅降低单设备的显存占用，缺点是需要频繁的all-reduce同步，对互联带宽要求极高，通常仅在NVLink互联的节点内部使用。

\textbf{流水线并行。}流水线并行（Pipeline Parallelism, PP）\cite{gpipe,huang2019gpipe}将模型按层切分为多个阶段，每个设备负责一个阶段的计算。GPipe\cite{gpipe}提出了micro-batch技术，将一个batch划分为多个micro-batch流水执行，以提高设备利用率。然而，GPipe需要完成所有micro-batch的前向传播后才开始反向传播，导致较高的内存峰值。

1F1B（One Forward One Backward）\cite{harlap2018pipedream}策略改进了这一问题，在micro-batch完成前向后立即启动其反向传播，降低了峰值内存。Chimera\cite{li2021chimera}提出了双向流水线，结合正向和逆向两个流水线进一步降低气泡率。Zero-Bubble\cite{qi2024zerobubble}创新性地将反向传播拆分为B-pass（计算激活梯度）和W-pass（计算权重梯度），通过灵活调度W-pass实现了接近零气泡的理想状态。

\textbf{序列并行。}序列并行（Sequence Parallelism, SP）\cite{li2022sequence}针对长序列场景设计，将激活值沿序列维度切分。Ring Attention\cite{liu2023ring}通过环状通信模式实现任意长度序列的注意力计算。LoongTrain\cite{gu2024loongtrain}结合SP与混合并行策略，优化了长上下文训练效率。

上述并行技术各有优劣，通常需要组合使用。Megatron-LM提出的3D并行（DP+TP+PP）已成为行业标准\cite{narayanan2021efficient}。然而，在长序列场景下，传统PP面临严峻挑战：激活值的规模与序列长度成正比，当序列足够长时，激活值通信量可能超过权重，此时"传递激活值"的范式不再高效。这正是本文WeiPipe方法所要解决的核心问题。

\subsection{长上下文训练技术}

长上下文能力是大模型的重要发展方向，但训练长序列模型面临独特挑战。

\textbf{注意力机制优化。}标准注意力的计算复杂度为$O(n^2)$，随序列长度平方增长。FlashAttention\cite{dao2023flashattention,dao2023flashattention2}通过分块计算和内存访问优化，将复杂度降低到$O(n)$的I/O操作。Ring Attention\cite{liu2023ring}支持超长序列的分布式注意力计算。LongLoRA\cite{chen2023longlora}通过LoRA微调扩展上下文长度。

\textbf{激活值管理。}长序列带来巨大的激活值存储压力。梯度检查点（Gradient Checkpointing）\cite{chen2016training}通过重计算策略用计算换存储，但增加了反向传播的计算量。混合精度训练\cite{micikevicius2018mixed}使用FP16/BF16存储激活值，降低存储开销。

这些技术虽然降低了单设备的内存压力，但在分布式训练中可能增加通信量。例如，更大的batch size和更长的序列长度都意味着更多的激活值需要在设备间传递。如何在长序列场景下重新设计通信模式，是本文WeiPipe工作的核心创新点。

\section{自动并行化方法研究}

面对复杂的硬件环境和多样的模型架构，手工设计并行策略变得日益困难。自动并行化（Automatic Parallelization）旨在自动发现最优的并行策略。本节从策略空间构建和搜索算法两个角度综述相关工作。

\subsection{基于预定义范式的自动搜索}

许多工作在预定义的并行范式框架内进行参数优化。

\textbf{3D并行参数调优。}Megatron-LM\cite{shoeybi2019megatron,narayanan2021efficient}定义了DP、TP、PP三维并行空间，自动搜索最优的并行度配置。这类方法搜索空间相对有限，仅考虑预定义范式的组合。

\textbf{分层搜索策略。}Alpa\cite{zheng2022alpa}采用两阶段搜索：第一阶段确定流水线划分，第二阶段优化每个stage内的张量并行。Piper\cite{tarnawski2021piper}提出类似的分层方法，但更关注流水线调度优化。这些方法将搜索问题分解，降低了复杂度，但也限制了搜索空间。

\textbf{约束优化方法。}Unity\cite{unger2022unity}将并行策略搜索建模为约束优化问题，使用整数线性规划求解。FlexFlow\cite{jia2019flexflow}定义了SOAP搜索空间（Sample-Operator-Attribute-Parameter），使用MCMC采样搜索。Tofu\cite{zhuang2023optimizing}专注于跨mesh的数据重分片优化。Galvatron\cite{miao2023galvatron}提出了混合并行性（Hybrid Parallelism）的自动搜索框架，将数据并行、张量并行、流水线并行和分片并行的组合纳入统一的动态规划搜索空间，通过精确的性能代价模型在数分钟内找到给定硬件配置下的最优并行策略。相比固定3D并行，Galvatron在超大规模模型上可实现显著的吞吐量提升。

上述方法的共同局限在于搜索空间由预定义范式限定，难以发现超越现有范式的新型策略。例如，它们无法发现ZeRO-3.5或WeiPipe这类需要精细控制数据流轨迹的策略。

\subsection{基于原语的策略空间构建}

为了突破预定义范式的局限，近期工作开始探索更灵活的抽象体系。

\textbf{nnScaler。}nnScaler\cite{lin2024nnscaler}提出了算子中心的原语系统，包括\texttt{op-trans}（算子变换）、\texttt{op-assign}（算子分配）和\texttt{op-order}（算子排序）。通过组合这些原语，可以表达超出3D并行的策略，如AlphaFold2所需的3F1B调度。然而，nnScaler的抽象仍以算子为中心，数据移动被视为算子调度的副作用，难以表达需要精细控制数据流轨迹的策略。

\textbf{Tessel。}Tessel\cite{lin2024tessel}专注于流水线调度的细粒度优化，在给定算子放置的情况下优化micro-batch调度顺序。它可以发现比传统1F1B更优的调度方案，但前提是算子放置已确定。

这些工作展示了原语化抽象的优势，但仍面临表达能力的根本局限：以算子为中心的视角难以显式控制数据的时空轨迹和生命周期。这正是本文CATFlow框架所要突破的核心障碍。

\subsection{数据流驱动的优化方法}

部分工作开始关注数据流在并行优化中的作用。

\textbf{数据移动感知优化。}GSPMD\cite{xu2021gspmd}将张量分片语义引入编译器，自动推导跨设备的数据移动。然而，它主要关注单算子内的切分，对跨算子的复杂数据流优化支持有限。

\textbf{内存优化策略。}ZeRO系列\cite{rajbhandari2020zero}和Activation Checkpointing\cite{chen2016training}本质上是对数据生命周期的管理，但它们是作为独立优化技术被设计，缺乏统一的抽象框架。

\textbf{重计算与卸载。}BPipe\cite{kim2023bpipe}结合流水线并行与激活重计算，通过精细的激活管理降低内存峰值。ZeRO-Offload\cite{ren2021zero}将优化器状态卸载到CPU，通过异步传输隐藏延迟。这些策略涉及复杂的数据流协调，在现有框架中难以自动发现。

本文的CATFlow框架提出了"数据一元论"的新视角，将数据作为第一性实体，算子计算被视为数据约束的自然结果。这一抽象使得表达和发现复杂的数据流优化策略成为可能。

\section{本章小结}

本章从大模型推理系统、分布式训练并行技术和自动并行化方法三个维度综述了相关研究工作，总结如下：

在推理能效优化方面，现有工作主要采用粗粒度的频率控制策略，未能充分利用推理负载的动态特性和阶段异质性。Prefill-Decode分离架构为精细化能效控制提供了新机会，但尚未被充分探索。

在训练通信优化方面，传统流水线并行的"传递激活值"范式在长序列场景下面临通信瓶颈。现有技术（如ZeRO、FlashAttention）主要关注内存优化，对通信模式的根本性重构研究不足。

在自动并行化方面，现有框架以算子为中心构建策略空间，表达能力受限，难以发现需要精细控制数据流轨迹的新型策略。如何构建更完备的抽象体系，使自动搜索有机会发现人类专家未曾设想的优化方案，是亟待解决的问题。

本文的三项核心工作——GreenServe、WeiPipe和CATFlow——正是针对上述局限性而设计，分别从能效优化、通信优化和策略发现三个维度推进大模型资源效率优化的研究前沿。
